%% toc: P15 171 171.11
\section{Generalization failure under high capacity or distribution shift: complete proofs}
\label{sec:annex-p15-w02b-p15-066}

Notation is that of \cref{def:p15-w02b-p15-066-two-laws}. For
\(z\in\mathcal Z\), \(\delta_z(A)=\mathbf1_A(z)\) is
the point mass; it is a probability measure on any sigma-field, whether
or not \(\{z\}\) is measurable.

Two facts about integrals are used repeatedly.
\emph{(I1) Dyadic approximation.} For measurable \(f:\mathcal Z\to[0,1]\)
put \(f_j=2^{-j}\lfloor2^jf\rfloor\). Each \(f_j\) is simple, measurable,
takes values in \(\{0,2^{-j},\ldots,1\}\subseteq[0,1]\), satisfies
\(f_j\leq f_{j+1}\) because \(2\lfloor a\rfloor\leq\lfloor2a\rfloor\),
and \(0\leq f-f_j<2^{-j}\), so \(f_j\uparrow f\) pointwise. By
\cref{thm:p15-monotone-integral}, \(\int f_j\,d\mu\uparrow\int f\,d\mu\)
for every measure \(\mu\).
\emph{(I2) Mixtures and point masses.} For measures
\(\mu,\nu\), constants \(a,b\geq0\) and measurable \(f\geq0\),
\(\int f\,d(a\mu+b\nu)=a\int f\,d\mu+b\int f\,d\nu\), and
\(\int f\,d\delta_z=f(z)\). For an indicator both identities are the
definitions; they extend to nonnegative simple functions by linearity of
the simple integral and to measurable \(f\geq0\) by approximating
\(f\) from below by simple functions and applying
\cref{thm:p15-monotone-integral} to each measure, which is (I1) when
\(f\leq1\) and the same construction truncated at level \(j\) in general.

\begin{annexitem}{annex:p15-w02b-p15-066-enclosing}{thm:p15-w02b-p15-066-enclosing}{Counting the classes that actually contain the outputs}
\begin{proof}
Item~1. By \cref{thm:p15-conc-selection-example} the probability that
some predictor has empirical risk zero is \(1-(1-2^{-n})^K\). Since
\(1-x\leq e^{-x}\) for real \(x\) (the convex function \(e^{-x}\) lies
above its tangent line \(1-x\) at \(0\)),
\[
 (1-2^{-n})^K\leq e^{-K2^{-n}}\leq e^{-\ln(1/\delta)}=\delta
\]
when \(K\geq2^n\ln(1/\delta)\). On the complementary event, of
probability at least \(1-\delta\), the minimum empirical risk is zero,
so every exact minimizer has \(\widehat R_n(\widehat h)=0\), while every
predictor, and so \(\widehat h\), has risk \(1/2\). If \(K\geq2^n\),
then \(\ln(2K/\delta)>\ln K\geq n\ln2\) because \(2/\delta>1\), so
\[
 b_n=\sqrt{\frac{\ln(2K/\delta)}{2n}}>\sqrt{\frac{\ln2}2}>\frac12,
\]
the last step because \(\ln2>1/2\), which holds since \(e<4\). If
\(\delta\) is fixed and \((\ln K)/n\to0\), then
\(\ln(2K/\delta)/(2n)=(\ln K+\ln(2/\delta))/(2n)\to0\).

Item~2. The class \(\{0,1\}^T\) of all binary functions on an
\(N\)-point domain \(T\) realizes every pattern on \(T\), so it shatters
\(T\); no larger set is contained in \(T\), so the dimension is \(N\).
\Cref{thm:p15-w02b-p15-065-lower-realizable} needs only a shattered set
of distinct points and measurable singletons, which the finite discrete
domain has.

Item~3. Every output has the form \(h_S=\mathbf1_{\mathcal X\setminus O(S)}\)
with \(O(S)\) the set of observed inputs, and \(1\leq|O(S)|\leq m\)
because the sample is nonempty and has \(m\) records. Conversely every
\(O\) with \(1\leq|O|\leq m\) occurs: a sample whose records run through
the points of \(O\) and then repeat one of them has positive probability
under the uniform law. Thus any fixed class containing \(f\) and all
possible outputs contains \(H_m\), and \(H_m\) is such a class.

Size. Distinct sets \(O\) give distinct functions, so \(H_m\) contains at
least \(\binom Nm\) functions \(\mathbf1_{\mathcal X\setminus O}\) with
\(|O|=m\). Writing
\(\binom Nm=\prod_{i=0}^{m-1}\frac{N-i}{m-i}\) and noting
\(\frac{N-i}{m-i}\geq\frac Nm\), which is equivalent to \(N\geq m\)
after cross-multiplication, gives \(r\geq\binom Nm\geq(N/m)^m\).

Dimension. Let \(N\geq m+2\). Take \(T\subseteq\mathcal X\) with
\(|T|=m+1\) and \(A\subseteq T\). For \(A=\varnothing\) use \(f\). For
nonempty \(A\ne T\) put \(O=T\setminus A\), so \(1\leq|O|\leq m\) and
\((\mathcal X\setminus O)\cap T=A\). For \(A=T\) choose a point
\(o\notin T\), which exists since \(N\geq m+2\), and put \(O=\{o\}\).
Thus \(T\) is shattered. If \(|T|=m+2\), a singleton \(A=\{a\}\subseteq T\)
is not realized: \(f\) gives the empty set, and
\((\mathcal X\setminus O)\cap T=\{a\}\) would need
\(O\supseteq T\setminus\{a\}\), a set of \(m+1\) points. No larger set is
shattered either, since its subsets of size \(m+2\) would be shattered.
The dimension is therefore \(m+1\).

The sample-size condition. \Cref{thm:p15-occam-finite} with \(r=|H_m|\)
requires \(m\geq(\ln r+\ln(1/\delta))/q\), \(q=-\ln(1-\varepsilon)\).
Since \(\ln(1/\delta)>0\) and, for \(N\geq2m\),
\(\ln r\geq m\ln(N/m)\geq m\ln2\), the condition gives
\(mq>m\ln2\), so \(1-\varepsilon<1/2\) and \(\varepsilon>1/2\). The
actual error bound \(1-m/N\geq1/2\) for \(N\geq2m\) is
\cref{thm:p15-occam-boundaries}, item~3.

Item~4. The memorizer of \cref{thm:p15-vcl-examples}, item~3, returns
the training label at a training point and one elsewhere. The training
label at \(X_i\) is \(\mathbf1_{\{X_i\geq1/2\}}\), so the returned
function is zero exactly on the finite set \(F\) of training points below
\(1/2\), which is \(\mathbf1_{[0,1]\setminus F}\). Each such function is
Borel because finite sets are closed. For a finite
\(T\subseteq[0,1/2)\) and \(A\subseteq T\), the choice \(F=T\setminus A\)
gives \(([0,1]\setminus F)\cap T=A\). Hence every finite subset of
\([0,1/2)\) is shattered, which has subsets of every finite size, and
the dimension is infinite.

Item~5. For finite \(T\subseteq\mathbb N\) and \(A\subseteq T\), the
finite set \(A\) itself satisfies \(A\cap T=A\), so every finite set is
shattered by the indicators of finite sets; the dimension is infinite.
In items~4 and~5 the domains \([0,1]\) with its Borel sets and
\(\mathbb N\) with all subsets have measurable singletons, so
\cref{thm:p15-w02b-p15-065-lower-infinite} applies to both classes.
\end{proof}
\end{annexitem}

\begin{annexitem}{annex:p15-w02b-p15-066-tv-transport}{thm:p15-w02b-p15-066-tv-transport}{The Hahn set of the difference and three consequences}
\begin{proof}
Item~1. By \cref{thm:p15-signed-measure-examples},
\(\alpha=Q-P\) is a signed measure; it is finite, and
\(\alpha(\mathcal Z)=0\). By \cref{thm:p15-hahn-decomposition} there is a
partition \(\mathcal Z=A\mathbin{\dot\cup}A^c\) with \(A\) positive and
\(A^c\) negative for \(\alpha\).

First, \(d=\alpha(A)\). For any \(E\in\mathcal B\),
\[
 \begin{aligned}
 Q(E)-P(E)&=\alpha(E\cap A)+\alpha(E\cap A^c)\leq\alpha(E\cap A)\\
 &\leq\alpha(E\cap A)+\alpha(A\setminus E)=\alpha(A),
 \end{aligned}
\]
using that \(E\cap A^c\subseteq A^c\) is negative and
\(A\setminus E\subseteq A\) is positive. Likewise
\[
 P(E)-Q(E)=-\alpha(E\cap A)-\alpha(E\cap A^c)\leq-\alpha(E\cap A^c)
 \leq-\alpha(A^c)=\alpha(A),
\]
using \(-\alpha(E\cap A^c)=-\alpha(A^c)+\alpha(A^c\setminus E)\leq-\alpha(A^c)\)
and \(\alpha(A^c)=-\alpha(A)\). So \(|P(E)-Q(E)|\leq\alpha(A)\) for all
\(E\), hence \(d\leq\alpha(A)\); and \(\alpha(A)=|Q(A)-P(A)|\leq d\).
Thus \(Q(A)-P(A)=d\).

Next, let \(f\) be simple with values in \([0,1]\), written
\(f=\sum_{k=1}^Kc_k\mathbf1_{E_k}\) with disjoint \(E_k\in\mathcal B\) and
\(c_k\in[0,1]\). Then
\[
 \int f\,dQ-\int f\,dP=\sum_kc_k\alpha(E_k)
 =\sum_kc_k\alpha(E_k\cap A)+\sum_kc_k\alpha(E_k\cap A^c).
\]
The terms \(\alpha(E_k\cap A)\) are nonnegative and the terms
\(\alpha(E_k\cap A^c)\) nonpositive. With \(0\leq c_k\leq1\) and
additivity,
\[
 \alpha(A^c)\leq\sum_k\alpha(E_k\cap A^c)
 \leq\int f\,dQ-\int f\,dP
 \leq\sum_k\alpha(E_k\cap A)\leq\alpha(A),
\]
where the outer inequalities hold because
\(\bigcup_kE_k\cap A\subseteq A\) and the remainder of \(A\) has
nonnegative measure, and symmetrically for \(A^c\). Since
\(\alpha(A^c)=-d\) and \(\alpha(A)=d\), the bound
\(|\int f\,dQ-\int f\,dP|\leq d\) holds for simple \(f\). For a general
measurable \(f\) with values in \([0,1]\), apply it to the dyadic
\(f_j\) of (I1) and let \(j\to\infty\); both integrals converge, and the
inequality is preserved. Taking \(f=\ell(h,\cdot)\) gives the risk
statement, and \(f=\mathbf1_A\) gives equality.

Item~2. On the event of \cref{thm:p15-conc-finite-learning}, every
\(h\) has \(|\widehat R_n(h)-R_P(h)|\leq b_n\); item~1 adds
\(|R_P(h)-R_Q(h)|\leq d\), and the triangle inequality gives the first
claim. For an \(\eta\)-approximate minimizer \(\widehat h\) the same
event gives \(R_P(\widehat h)\leq\min_hR_P(h)+2b_n+\eta\). Let
\(h^\ast\) minimize \(R_Q\) over the finite class. Then
\[
 R_Q(\widehat h)\leq R_P(\widehat h)+d
 \leq R_P(h^\ast)+2b_n+\eta+d
 \leq R_Q(h^\ast)+2b_n+2d+\eta .
\]
The first and last steps use item~1 for the fixed candidates
\(\widehat h(\omega)\) and \(h^\ast\); item~1 is a deterministic
inequality, so it applies to whichever candidate the sample selects.
In the setting of \cref{thm:p15-vcl-erm}, on an event where
\(\sup_h|\widehat R_m(h)-R_P(h)|\leq u\), that theorem gives
\(R_P(\widehat h)\leq\inf_hR_P(h)+2u+\eta\). For every \(h\),
\(\inf_gR_P(g)\leq R_P(h)\leq R_Q(h)+d\), so
\(\inf_hR_P(h)\leq\inf_hR_Q(h)+d\), and
\(R_Q(\widehat h)\leq R_P(\widehat h)+d\) completes the bound. In the
realizable setting, \(R_Q(h,f)=Q\{h\ne f\}\) and
\(R_P(h,f)=P\{h\ne f\}\) for the measurable set \(\{h\ne f\}\), so
\(R_Q(h,f)\leq R_P(h,f)+d\leq\varepsilon+d\) for every consistent
\(h\) on the event of \cref{thm:p15-occam-finite}.

Item~3. Let \(\mathcal Y\) be finite with the discrete sigma-field. For
measurable \(C\subseteq\mathcal X\times\mathcal Y\) and \(y\in\mathcal Y\),
the section \(C_y=\{x:(x,y)\in C\}\) is measurable: the sets with
measurable sections form a sigma-field, since sections commute with
complements and countable unions, and it contains every rectangle
\(B\times D\), whose section is \(B\) or \(\varnothing\). Hence
\(g_C(x)=\sum_y\mathbf1_{C_y}(x)p(y\mid x)\) is measurable with values in
\([0,1]\), and by the covariate-shift hypothesis
\(Q(C)-P(C)=\int g_C\,dQ_{\mathcal X}-\int g_C\,dP_{\mathcal X}\).
Item~1 on \(\mathcal X\) gives
\(|Q(C)-P(C)|\leq\|P_{\mathcal X}-Q_{\mathcal X}\|_{\mathrm{TV}}\), so
\(\|P-Q\|_{\mathrm{TV}}\leq\|P_{\mathcal X}-Q_{\mathcal X}\|_{\mathrm{TV}}\).
For \(C=B\times\mathcal Y\), \(g_C=\mathbf1_B\) and
\(Q(C)-P(C)=Q_{\mathcal X}(B)-P_{\mathcal X}(B)\); the supremum over \(B\)
gives the reverse inequality.
\end{proof}
\end{annexitem}

\begin{annexitem}{annex:p15-w02b-p15-066-density-transport}{thm:p15-w02b-p15-066-density-transport}{From indicators to all nonnegative functions}
\begin{proof}
First, \(1=Q(\mathcal Z)=\int w\,dP\leq\int B\,dP=B\) by monotonicity in
\cref{thm:p15-monotone-integral}.

For \(f=\mathbf1_E\) the identity \(\int f\,dQ=\int fw\,dP\) is the
hypothesis \(Q(E)=\int_Ew\,dP\). For a nonnegative simple
\(f=\sum_kc_k\mathbf1_{E_k}\), both sides are linear in \(f\) by
\cref{thm:p15-nonnegative-integral-laws}, so the identity holds. For
measurable \(f\geq0\), choose nonnegative simple \(f_j\uparrow f\)
(for instance \(f_j=\min\{j,2^{-j}\lfloor2^jf\rfloor\}\), with
\(f_j=j\) where \(f=\infty\)); then \(f_jw\uparrow fw\), and
\cref{thm:p15-monotone-integral}, applied under \(Q\) and under \(P\),
passes the identity to the limit. Since \(fw\leq Bf\), monotonicity
and \cref{thm:p15-nonnegative-integral-laws} give
\(\int fw\,dP\leq B\int f\,dP\).

Taking \(f=\ell(h,\cdot)\) gives \(R_Q(h)\leq BR_P(h)\). In the
realizable setting, \(f=\mathbf1_{\{h\ne f_0\}}\) for the target \(f_0\)
gives \(R_Q(h,f_0)\leq BR_P(h,f_0)\leq B\varepsilon\) for every
consistent \(h\) on the event of \cref{thm:p15-occam-finite}.

Sharpness. Let \(0<\varepsilon\leq1/B\). If \(\varepsilon=1\), then
\(B=1\) and \(Q=P=\delta_1\), with \(w=1\). Otherwise put
\(w(1)=B\) and \(w(0)=(1-B\varepsilon)/(1-\varepsilon)\), which lies in
\([0,1]\subseteq[0,B]\) because \(\varepsilon\leq B\varepsilon\leq1\).
Then \(\int_{\{1\}}w\,dP=B\varepsilon=Q\{1\}\) and
\(\int_{\{0\}}w\,dP=1-B\varepsilon=Q\{0\}\), so \(w\) is a density of
\(Q\) relative to \(P\) bounded by \(B\). With target \(0\) and
\(h(x)=x\), \(R_P(h,0)=P\{1\}=\varepsilon\) and
\(R_Q(h,0)=Q\{1\}=B\varepsilon\).

Finally, if \(P(A)=0\), then \(Q(A)=\int_Aw\,dP\leq BP(A)=0\), so no
bounded density exists when \(Q(A)>0=P(A)\).
\end{proof}
\end{annexitem}

\begin{annexitem}{annex:p15-w02b-p15-066-shift-attained}{thm:p15-w02b-p15-066-shift-attained}{Four shifts}
\begin{proof}
Item~1. On \(\{0,1\}\) the events are \(\varnothing,\{0\},\{1\},\{0,1\}\),
and \(|Q(E)-P(E)|\) equals \(d\) for \(\{0\}\) and \(\{1\}\) and zero for
the others, so \(\|P-Q\|_{\mathrm{TV}}=d\). Under \(P=\delta_0\) every
training input is \(0\), where \(h(0)=0=f(0)\); so \(h\) is consistent
with every sample, and \(R_P(h,f)=P\{1\}=0\). Under \(Q\),
\(R_Q(h,f)=Q\{h\ne f\}=Q\{1\}=d\). For \(d=1\) these are the laws of
\cref{thm:p15-occam-boundaries}, item~2.

Item~2. The laws are laws of the label alone; \(P\{Y=1\}=1/2+\gamma\in(1/2,1]\)
because \(0<\gamma\leq1/2\), and \(Q\{Y=1\}=1/2+\gamma-d\in[0,1/2)\)
because \(\gamma<d\leq1/2+\gamma\). On \(\{0,1\}\) the total variation is
\(|P\{Y=1\}-Q\{Y=1\}|=d\), as in item~1. The risks are
\[
 \begin{aligned}
 R_P(h_1)&=\tfrac12-\gamma,&\quad R_P(h_2)&=\tfrac12+\gamma,\\
 R_Q(h_1)&=\tfrac12-\gamma+d,&\quad R_Q(h_2)&=\tfrac12+\gamma-d .
 \end{aligned}
\]
Since \(\gamma>0\), \(h_1\) is the unique minimizer of \(R_P\). Since
\(d>\gamma\), \(R_Q(h_2)<R_Q(h_1)\), and the difference is
\(2(d-\gamma)\). With \(W_i=\mathbf1\{Y_i=0\}\),
\(\widehat R_n(h_1)=\bar W\) and \(\widehat R_n(h_2)=1-\bar W\); with
ties resolved in favour of \(h_1\), the minimizer returns \(h_2\) only
if \(\bar W>1/2\). The \(W_i\) are independent, lie in \([0,1]\) and
have mean \(1/2-\gamma\). By \cref{thm:p15-conc-unequal-ranges} with
\(V_n=n\) and \(t=n\gamma\),
\[
 P\{\bar W>\tfrac12\}\leq
 P\Bigl\{\sum_i\bigl(W_i-(\tfrac12-\gamma)\bigr)\geq n\gamma\Bigr\}
 \leq e^{-2(n\gamma)^2/n}=e^{-2n\gamma^2}.
\]

Item~3. For \(E\in\mathcal B\),
\(Q_d(E)-P(E)=d(\mathbf1_E(z_1)-P(E))\), whose absolute value is at most
\(d\); so \(\|P-Q_d\|_{\mathrm{TV}}\leq d\). By (I2),
\(R_{Q_d}(h)=(1-d)R_P(h)+d\,\ell(h,z_1)=(1-d)R_P(h)+d\). The training
sample has law \(P^{\otimes n}\) whatever \(Q\) is, so a statistic
\(T\) valid for every \(Q\) in the ball is valid for \(Q_d\):
\(P^{\otimes n}\{(1-d)R_P(h)+d\leq T\}\geq1-\delta\). If the
requirement holds for every law, take \(Q=Q_1=\delta_{z_1}\), for which
\(R_Q(h)=1\).

Item~4. The set \(\{t\}\) is Borel, and
\(Q\{t\}-P\{t\}=1\), so \(\|P-Q\|_{\mathrm{TV}}=1\), the largest
possible value. With label \(0\), \(R_P(h_s)=h_s(0)=\mathbf1_{\{s\leq0\}}\)
and \(R_Q(h_s)=h_s(t)=\mathbf1_{\{s\leq t\}}\) by (I2). Their difference
is one exactly when \(0<s\leq t\) and zero otherwise.
\end{proof}
\end{annexitem}

\begin{annexitem}{annex:p15-w02b-p15-066-worked}{thm:p15-w02b-p15-066-worked}{The worked numbers}
\begin{proof}
Item~1. Since \(\ln20<2.996\) and \(2^{10}\cdot2.996=3067.904<3068\),
\(K\geq2^{10}\ln20\). By \cref{thm:p15-w02b-p15-066-enclosing}, item~1,
with \(\delta=1/20\), the zero-fit event has probability at least
\(19/20\), and on it the minimizer has training error zero and risk
\(1/2\). Here \(2K/\delta=2\cdot3068\cdot20=122720\), and
\(K\geq2^{10}\) gives \(b_{10}>1/2\) by the same item.

Item~2. For \(n=1016\), \(r=4\) and \(\delta=1/20\),
\(b_n=\sqrt{\ln160/2032}\). Since \(\ln160<5.08\) and
\(5.08/2032=0.0025\), \(b_n<\sqrt{0.0025}=1/20\), so \(2b_n<1/10\). By
\cref{thm:p15-conc-noisy-example} the risks of \(0,1,x,1-x\) under \(P\)
are \(1/2,1/2,1/4,3/4\). On the event of
\cref{thm:p15-conc-finite-learning}, the minimizer has
\(R_P(\widehat h)\leq1/4+2b_n<1/2\), so \(\widehat h=x\).

(a) Under \(Q\), \(Q\{X=0\}=1-q\), \(Q\{X=1\}=q\), and the conditional
label probabilities are those of \(P\). The constant \(0\) errs when
\(Y=1\): \((1-q)/4+3q/4=1/4+q/2\). The constant \(1\) errs when
\(Y=0\): \(3(1-q)/4+q/4=3/4-q/2\). The predictor \(x\) errs when
\(Y\ne X\): \((1-q)/4+q/4=1/4\). The predictor \(1-x\) errs when
\(Y=X\): \(3(1-q)/4+3q/4=3/4\). Since \(0\leq q\leq1\), the value \(1/4\)
is the least of the four. By \cref{thm:p15-w02b-p15-066-tv-transport},
item~3, \(\|P-Q\|_{\mathrm{TV}}\) is the distance between the input laws
on \(\{0,1\}\), which is \(|q-1/2|\) as in the proof of
\cref{thm:p15-w02b-p15-066-shift-attained}, item~1. The transfer bound
of \cref{thm:p15-w02b-p15-066-tv-transport}, item~2, with \(\eta=0\)
reads \(R_Q(x)\leq1/4+2b_n+2|q-1/2|\leq1/4+1/10+2|q-1/2|\).

(b) The four records have probabilities
\[
 \begin{array}{c|cccc}
 (x,y)&(0,0)&(0,1)&(1,0)&(1,1)\\\hline
 P&3/8&1/8&1/8&3/8\\
 Q'&1/8&3/8&3/8&1/8
 \end{array}
\]
On a finite space with all singletons measurable, the event
\(A=\{(0,1),(1,0)\}\) where \(Q'\) exceeds \(P\) gives
\(Q'(A)-P(A)=1/2\); every event \(E\) has
\(Q'(E)-P(E)=Q'(E\cap A)-P(E\cap A)+Q'(E\setminus A)-P(E\setminus A)\leq1/2\)
and, symmetrically, \(P(E)-Q'(E)\leq P(A^c)-Q'(A^c)=1/2\). Hence
\(\|P-Q'\|_{\mathrm{TV}}=1/2\). Under \(Q'\) the constant \(0\) errs on
\(\{(0,1),(1,1)\}\), probability \(1/2\); the constant \(1\) on
\(\{(0,0),(1,0)\}\), probability \(1/2\); \(x\) on
\(\{(0,1),(1,0)\}\), probability \(3/4\); and \(1-x\) on
\(\{(0,0),(1,1)\}\), probability \(1/4\). So
\(R_{Q'}(x)-R_P(x)=3/4-1/4=1/2=\|P-Q'\|_{\mathrm{TV}}\), and the excess
of \(x\) over the best predictor under \(Q'\) is \(3/4-1/4=1/2\).
\end{proof}
\end{annexitem}
